% ZNS Whitepaper diagrams
% Include individual figures via \input{figures/<name>}

% === Architecture Stack ===
% Shows the four actors and the trust boundary between each.
% Place in §02 (System and Trust Model)

\newcommand{\znsarchitecture}{%
\begin{tikzpicture}[
  node distance=1.4cm,
  box/.style={draw, rounded corners=2pt, minimum width=2.8cm, minimum height=1.1cm, align=center, font=\small},
  arrow/.style={-Stealth, thick},
  label/.style={font=\scriptsize\itshape, text=gray}
]

% Actors
\node[box, fill=gray!8]  (wallet)  {Wallet};
\node[box, fill=gray!8, below=of wallet]  (mint)    {Mint (TEE)};
\node[box, fill=gray!15, below=of mint]    (chain)   {Zcash chain};
\node[box, fill=gray!8, below=of chain]  (resolver) {Resolver};

% Arrows between actors
\draw[arrow] (wallet) -- node[right, label] {request} (mint);
\draw[arrow] (mint) -- node[right, label] {self-send} (chain);
\draw[arrow] (chain) -- node[right, label] {scan} (resolver);
\draw[arrow, dashed] (resolver.west) to[bend left=40] node[left, label] {resolve} (wallet.west);

% Trust boundary annotations
\node[font=\scriptsize, text=gray, right=0.3cm of mint.east, anchor=west] (t1) {attestation};
\node[font=\scriptsize, text=gray, right=0.3cm of chain.east, anchor=west] (t2) {consensus};
\node[font=\scriptsize, text=gray, right=0.3cm of resolver.east, anchor=west] (t3) {deterministic};

\end{tikzpicture}%
}

% === Name Lifecycle State Machine ===
% Shows the four states and transitions.
% Place in §05 (Canonical History and State Derivation)

\newcommand{\znslifecycle}{%
\begin{tikzpicture}[
  node distance=2.8cm,
  state/.style={draw, circle, minimum size=1.8cm, align=center, font=\small},
  arrow/.style={-Stealth, thick},
  tlabel/.style={font=\scriptsize, fill=white, inner sep=2pt}
]

\node[state, fill=gray!8] (unclaimed) {un-\\claimed};
\node[state, fill=gray!12, right=of unclaimed] (live) {live};
\node[state, fill=gray!8, right=of live] (released) {re-\\leased};

% Transitions
\draw[arrow] (unclaimed) -- node[tlabel, above] {\action{claim}} (live);
\draw[arrow] (live.north) to[bend left=40] node[tlabel, above] {\action{update}} (live.north);
\draw[arrow] (live) -- node[tlabel, above] {\action{release}} (released);
\draw[arrow] (released) to[bend left=35] node[tlabel, below] {\action{claim}} (live);

% Predecessor annotations
\node[font=\scriptsize\itshape, text=gray, below=0.1cm of unclaimed] {prev = zero};
\node[font=\scriptsize\itshape, text=gray, below=0.1cm of live] {prev = live tip};
\node[font=\scriptsize\itshape, text=gray, below=0.1cm of released] {prev = release tip};

\end{tikzpicture}%
}

% === Predecessor Chain ===
% Shows how prev_rcm links name notes into one continuous history.
% Place in §05 (Canonical History and State Derivation)

\newcommand{\znschain}{%
\begin{tikzpicture}[
  node distance=0.8cm,
  note/.style={draw, rounded corners=2pt, minimum width=2.2cm, minimum height=0.9cm, align=center, font=\small},
  arrow/.style={-Stealth, thick},
  rlabel/.style={font=\scriptsize, text=gray}
]

% Name notes
\node[note, fill=gray!8] (c1) {\action{claim}\\\scriptsize prev = $\bot$};
\node[note, fill=gray!12, right=of c1] (u1) {\action{update}\\\scriptsize prev = rcm$_1$};
\node[note, fill=gray!8, right=of u1] (r1) {\action{release}\\\scriptsize prev = rcm$_2$};
\node[note, fill=gray!12, right=of r1] (c2) {\action{claim}\\\scriptsize prev = rcm$_3$};

% Arrows
\draw[arrow] (c1) -- (u1);
\draw[arrow] (u1) -- (r1);
\draw[arrow] (r1) -- (c2);

% Derived rcm labels below
\node[rlabel, below=0.15cm of c1] {rcm$_1$};
\node[rlabel, below=0.15cm of u1] {rcm$_2$};
\node[rlabel, below=0.15cm of r1] {rcm$_3$};
\node[rlabel, below=0.15cm of c2] {rcm$_4$};

% Registration period bracket
\node[font=\scriptsize\itshape, text=gray, above=0.3cm of c1.north west, anchor=south west] {registration 1};
\draw[gray, thick] ([yshift=0.2cm]c1.north west) -- ([yshift=0.2cm]r1.north east);

\node[font=\scriptsize\itshape, text=gray, above=0.7cm of c2.north west, anchor=south west] {registration 2};
\draw[gray, thick] ([yshift=0.85cm]c2.north west) -- ([yshift=0.85cm]c2.north east);

\end{tikzpicture}%
}

\newcommand{\znsauthorizationflow}{%
\begin{tikzpicture}[
  node distance=1.6cm,
  actor/.style={draw, rounded corners=2pt, minimum width=2.6cm, minimum height=1.0cm, align=center, font=\small},
  chainbox/.style={draw, dashed, rounded corners=2pt, minimum width=2.6cm, minimum height=0.7cm, align=center, font=\small\itshape},
  arrow/.style={-Stealth, thick},
  label/.style={font=\scriptsize\itshape, text=gray, fill=white, inner sep=2pt}
]

\node[actor, fill=gray!8] (wallet) {Wallet};
\node[actor, fill=gray!8, right=3.2cm of wallet] (mint) {Mint (TEE)};
\node[chainbox, fill=gray!15, below=1.8cm of wallet, xshift=1.6cm] (chain) {Zcash chain};

\draw[arrow] (wallet) -- node[label, above] {request memo} (mint);
\draw[arrow] (mint) -- node[label, sloped, above] {OTP relay memo} (chain);
\draw[arrow] (chain) -- node[label, sloped, above] {OTP relay memo} (wallet);
\draw[arrow] (wallet) -- node[label, sloped, below] {request + OTP} (chain);
\draw[arrow] (chain) -- node[label, sloped, below] {request + OTP} (mint);
\draw[arrow, dashed] (mint) -- node[label, right] {Name Note self-send} (chain);

\node[font=\scriptsize, text=gray, below=0.15cm of wallet] {controller UA};
\node[font=\scriptsize, text=gray, below=0.15cm of mint] {OTP account};

\end{tikzpicture}%
}

% === ZNS vs ENS vs Namecoin Comparison ===
% Shows where each system places name state relative to consensus.
% Place in §02 (System and Trust Model) or a new comparison section.

\newcommand{\znscomparison}{%
\begin{tikzpicture}[
  node distance=0.6cm,
  sysbox/.style={draw, rounded corners=2pt, minimum width=3.2cm, minimum height=2.4cm, align=center, font=\small},
  consensus/.style={draw, dashed, minimum width=3.2cm, minimum height=0.8cm, align=center, font=\small\itshape, fill=gray!5},
  arrow/.style={-Stealth, thick}
]

% Namecoin: name state inside consensus
\node[consensus] (nc-cons) {consensus};
\node[sysbox, fill=gray!12, below=0.1cm of nc-cons, minimum height=0.6cm] (nc-name) {name state};
\node[font=\small\bfseries, above=0.1cm of nc-cons] {Namecoin};

% ENS: name state inside consensus (contract)
\node[consensus, right=1.5cm of nc-cons] (ens-cons) {consensus};
\node[sysbox, fill=gray!12, below=0.1cm of ens-cons, minimum height=0.6cm] (ens-name) {contract state};
\node[font=\small\bfseries, above=0.1cm of ens-cons] {ENS};

% ZNS: name state above consensus
\node[consensus, right=1.5cm of ens-cons] (zns-cons) {consensus};
\node[sysbox, fill=gray!8, below=0.5cm of zns-cons, minimum height=0.6cm] (zns-name) {Mint + Resolver};
\node[sysbox, fill=gray!12, below=0.1cm of zns-name, minimum height=0.4cm] (zns-tx) {Orchard tx};
\node[font=\small\bfseries, above=0.1cm of zns-cons] {ZNS};

\end{tikzpicture}%
}